\chapter{\texorpdfstring{附\quad 录}{附录}} %这里多出来的代码是用来消除hyperref包引入的警告的，\texorpdfstring{}{}命令的作用是告诉hyperref包在生成PDF书签时使用第二个参数，而不是第一个参数中的特殊字符。

\begin{appendixenv}

附录的有无，根据毕业设计（论文）情况而定，内容一般包括正文内不便列出的冗长公式推导、辅助性数学工具、符号说明（含缩写）、计算程序及说明等。附录另起一页。附录依序用大写正体A、B、C…编序号，如：附录A。附录中的图、表、公式、参考文献等另行编序号，与正文分开，也一律用阿拉伯数字编码，但在数码前冠以附录序码，如：图A01、表B-2、公式（B-3）、文献[A5]等。“附录”两字之间空一个全角空格或两个半角空格，采用不编号章标题样式；内容应采用正文样式。

\section{格式演示}

\figuremacro{htbp}
            {晶体坐标系和样品坐标系示意图.jpg}
            {晶体坐标系和样品坐标系示意图}
            {0.35}
            {晶体2}

\vspace{-20pt}  %单纯为了压缩到一页[doge]

\subsection{公式与表}
% \subsubsection{公式与表}

\vspace{-20pt}  %单纯为了压缩到一页[doge]

\begin{equation}
    \begin{cases}
        F_{si} &= k_{si}(z'_i-z_i)+c_{si}(\dot{z}'_i-\dot{z}_i) \\
        F_{ti} &= k_{ti}(z_i-y_i(x_i,t))+c_{ti}(\dot{z}_i-\dot{y}_i(x_i,t))
    \end{cases}
    \label{apequ}
\end{equation}

\vspace{-12pt}  %单纯为了压缩到一页[doge]

\tablemacro{htbp}
            {典型微尺度力学实验方法的基本信息}
            {ccc}
            {\songtibf{实验方法} & \songtibf{主要测量对象} & \songtibf{空间分辨率} \\}
            {原子力显微镜 & 表面力 & $0.1 nm$\cite{Zhang2002} \\
            透射电镜 & 晶格结构、位错 & $0.1 nm$\cite{Borko1978} \\
            X射线衍射 & 应变 & $1 \mu m$\cite{Hao2008,Dupont1974,Wang1997} \\
            同步辐射 & 内部三维结构与变形 & 约$100 nm$\cite{Liang2024,Koya1995} \\
            显微拉曼 & 应变 & $250 nm$\cite{Deng2006} \\}
            {基本信息2}

\end{appendixenv}

\clearpage
